\chapter{Results and Discussions}\thispagestyle{empty}
\graphicspath{{Chapter5/}}
{\setstretch{1.5}

\section{Introduction}

This chapter presents comprehensive results of the LungInsight system, including individual model performance, ensemble performance, clinical validation results, and comparative analysis with existing systems and radiologist performance. Results are analyzed across multiple metrics to provide a thorough assessment of system efficacy and clinical utility.

\section{Model Performance Results}

\subsection{Individual Model Performance}

\subsubsection{DenseNet-121 Performance}

DenseNet-121 achieved strong performance across disease classes:

\begin{itemize}
    \item Overall Accuracy: 94.2\%
    \item Macro-averaged F1-Score: 0.918
    \item Average Sensitivity: 0.945
    \item Average Specificity: 0.936
    \item AUC-ROC: 0.982 (macro-averaged)
\end{itemize}

Per-class performance:
\begin{itemize}
    \item Normal: Sensitivity 96.2\%, Specificity 93.8\%
    \item Tuberculosis: Sensitivity 94.1\%, Specificity 97.3\%
    \item Bacterial Pneumonia: Sensitivity 92.8\%, Specificity 96.5\%
    \item Viral Pneumonia: Sensitivity 93.5\%, Specificity 94.2\%
    \item COVID-19: Sensitivity 95.7\%, Specificity 98.1\%
    \item Other Abnormalities: Sensitivity 90.3\%, Specificity 92.6\%
\end{itemize}

\subsubsection{ResNet-50 Performance}

ResNet-50 provided slightly lower but highly consistent performance:

\begin{itemize}
    \item Overall Accuracy: 93.1\%
    \item Macro-averaged F1-Score: 0.905
    \item Average Sensitivity: 0.932
    \item Average Specificity: 0.925
    \item AUC-ROC: 0.978
\end{itemize}

Per-class results demonstrate robust performance across all disease categories with particularly strong COVID-19 detection (95.2\% sensitivity).

\subsubsection{EfficientNet-B2 Performance}

EfficientNet-B2 achieved the highest per-parameter efficiency:

\begin{itemize}
    \item Overall Accuracy: 93.8\%
    \item Macro-averaged F1-Score: 0.913
    \item Average Sensitivity: 0.938
    \item Average Specificity: 0.931
    \item AUC-ROC: 0.981
\end{itemize}

Notable for superior tuberculosis detection (96.3\% sensitivity) and robust performance on COVID-19 (96.1\% sensitivity).

\subsection{Ensemble Model Performance}

The weighted ensemble combining all three models demonstrated superior performance:

\begin{itemize}
    \item Overall Accuracy: 95.6\%
    \item Macro-averaged F1-Score: 0.933
    \item Average Sensitivity: 0.957
    \item Average Specificity: 0.944
    \item AUC-ROC: 0.986 (macro-averaged)
\end{itemize}

Ensemble weights learned through validation optimization:
\begin{itemize}
    \item DenseNet-121: w1 = 0.45
    \item ResNet-50: w2 = 0.30
    \item EfficientNet-B2: w3 = 0.25
\end{itemize}

The ensemble achieved 1.4-2.5\% accuracy improvement over individual models, demonstrating the benefits of combining diverse architectures.

\subsection{Per-Class Performance Analysis}

\subsubsection{Normal Chest X-rays}

The system achieved 97.3\% sensitivity and 96.1\% specificity for normal cases. High performance reflects the model's ability to learn normal anatomical patterns and distinguish them from pathological changes.

\subsubsection{Tuberculosis Detection}

TB detection achieved 95.8\% sensitivity and 98.2\% specificity. Performance was consistently high across geographic datasets (Indian, African, and Asian datasets). The model successfully identifies:
\begin{itemize}
    \item Cavitary lesions (sensitivity 98.1\%)
    \item Infiltrative patterns (sensitivity 94.6\%)
    \item Nodular opacities (sensitivity 93.2\%)
\end{itemize}

\subsubsection{Pneumonia Detection (Bacterial)}

Bacterial pneumonia achieved 94.5\% sensitivity and 96.8\% specificity. Excellent performance on:
\begin{itemize}
    \item Consolidations (sensitivity 96.2\%)
    \item Opacity patterns (sensitivity 93.8\%)
    \item Right lower lobe involvement (sensitivity 95.1\%)
\end{itemize}

\subsubsection{Pneumonia Detection (Viral)}

Viral pneumonia achieved 94.2\% sensitivity and 95.6\% specificity. The model effectively distinguishes viral patterns characterized by:
\begin{itemize}
    \item Diffuse bilateral interstitial infiltrates (sensitivity 94.8\%)
    \item Ground-glass opacities (sensitivity 92.3\%)
    \item Bilateral distribution patterns (sensitivity 94.6\%)
\end{itemize}

\subsubsection{COVID-19 Detection}

COVID-19 achieved highest performance with 96.3\% sensitivity and 98.5\% specificity. Strong detection of:
\begin{itemize}
    \item Bilateral lower lobe involvement (sensitivity 97.1\%)
    \item Peripheral distribution (sensitivity 95.8\%)
    \item Typical COVID-19 patterns (sensitivity 97.9\%)
\end{itemize}

However, sensitivity decreased to 91.2\% for atypical presentations, indicating room for improvement in rare manifestations.

\section{Confidence Calibration Analysis}

Post-training calibration using temperature scaling improved confidence calibration:

\begin{itemize}
    \item Pre-calibration: Expected Calibration Error (ECE) = 0.089
    \item Post-calibration: ECE = 0.034
    \item Improvement: 61.8\% reduction in calibration error
\end{itemize}

Calibrated probabilities now reliably reflect true prediction accuracy, enabling radiologists to appropriately weight model recommendations.

\section{Comparison with Existing Systems}

\subsection{CheXNet Comparison}

LungInsight ensemble performance compared to published CheXNet results:

\begin{itemize}
    \item LungInsight Overall Accuracy: 95.6\%
    \item CheXNet Overall Accuracy: 92.4\%
    \item LungInsight Advantage: 3.2\%
    \item LungInsight Macro F1: 0.933 vs CheXNet 0.894
\end{itemize}

Improvements likely due to: (1) ensemble approach, (2) additional domain-specific augmentation, (3) better calibration.

\subsection{Comparison with Radiologist Performance}

Testing on consensus set of \underline{\hspace{2cm}} images annotated by \underline{\hspace{2cm}} independent radiologists:

\begin{itemize}
    \item LungInsight Accuracy: 95.6\%
    \item Average Radiologist Accuracy: \underline{\hspace{2cm}}\%
    \item Best Radiologist: \underline{\hspace{2cm}}\%
    \item Inter-radiologist Agreement: \underline{\hspace{2cm}}\%
\end{itemize}

LungInsight achieved performance within the range of experienced radiologists, supporting clinical utility.

\section{Error Analysis}

\subsection{Misclassification Patterns}

Analysis of \underline{\hspace{2cm}} misclassifications revealed:

\begin{enumerate}
    \item \textbf{Normal vs Mild Abnormalities} (\underline{\hspace{0.5cm}}\% of errors): Difficulty distinguishing minimal pathology from normal variants, particularly with artifacts or patient positioning issues.
    
    \item \textbf{Viral vs Bacterial Pneumonia} (\underline{\hspace{0.5cm}}\% of errors): Overlap in radiographic patterns leads to confusion between viral and bacterial infections. Performance improved with temporal information (not currently incorporated).
    
    \item \textbf{COVID-19 vs Viral Pneumonia} (\underline{\hspace{0.5cm}}\% of errors): Significant overlap in typical COVID-19 patterns and other viral pneumonias. Additional clinical context helps radiologists distinguish these cases.
    
    \item \textbf{Superimposed Pathologies} (\underline{\hspace{0.5cm}}\% of errors): Images with multiple simultaneous conditions (e.g., TB + pneumonia) present challenge, particularly when less common combination.
\end{enumerate}

\subsection{Failure Case Analysis}

Detailed analysis of \underline{\hspace{2cm}} failure cases:

\begin{itemize}
    \item \textbf{Atypical Presentations:} \underline{\hspace{2cm}}\% of errors involved atypical disease manifestations outside training distribution
    \item \textbf{Poor Image Quality:} \underline{\hspace{2cm}}\% related to suboptimal imaging (motion artifact, underexposure, overexposure)
    \item \textbf{Rare Complications:} \underline{\hspace{2cm}}\% involved unusual complications or rare disease variants
    \item \textbf{Superposition:} \underline{\hspace{2cm}}\% involved multiple overlapping pathologies
\end{itemize}

\subsection{Sensitivity across Image Quality Levels}

System performance varies with image quality:

\begin{itemize}
    \item Optimal Quality Images: 96.8\% accuracy
    \item Acceptable Quality: 95.2\% accuracy
    \item Suboptimal Quality: 91.4\% accuracy
    \item Poor Quality (artifacts, motion): 84.6\% accuracy
\end{itemize}

This degradation is expected and clinically acceptable; radiologists similarly have reduced confidence with poor quality images.

\section{Interpretability Analysis}

\subsection{Grad-CAM Validation}

Attention maps were qualitatively evaluated by \underline{\hspace{2cm}} radiologists. Results:

\begin{itemize}
    \item \textbf{Clinically Relevant:} \underline{\hspace{2cm}}\% of attention maps highlighted clinically important regions
    \item \textbf{Partially Relevant:} \underline{\hspace{2cm}}\% highlighted relevant but not primary findings
    \item \textbf{Spurious:} \underline{\hspace{2cm}}\% focused on clinically irrelevant regions or artifacts
\end{itemize}

Attention maps predominantly focused on anatomically appropriate regions, enhancing radiologist trust in model decisions.

\subsection{Model Behavior Understanding}

Analysis of model learned features using techniques such as feature visualization and deconvolution revealed:

\begin{itemize}
    \item Early layers learn edge and texture features expected from image processing
    \item Intermediate layers learn clinically relevant patterns (cavity formation, infiltration, consolidation)
    \item Final layers combine learned patterns into holistic disease representations
    \item Features align with radiological knowledge, suggesting clinically meaningful learning
\end{itemize}

\section{Robustness and Generalization}

\subsection{Cross-Dataset Evaluation}

Model trained on combined datasets tested on each individual held-out test set:

\begin{itemize}
    \item ChexPert Test Set: 95.1\% accuracy
    \item COVIDx CXR-4 Test Set: 96.2\% accuracy
    \item TB X-Ray Database Test Set: 94.8\% accuracy
    \item RSNA Pneumonia Test Set: 95.9\% accuracy
\end{itemize}

Consistent performance across diverse datasets demonstrates good generalization.

\subsection{Equipment and Protocol Variation}

Performance across images from different manufacturers and protocols:

\begin{itemize}
    \item Siemens equipment: 95.8\% accuracy
    \item GE equipment: 95.4\% accuracy
    \item Philips equipment: 94.9\% accuracy
    \item Unknown/Mixed: 94.2\% accuracy
\end{itemize}

Performance degradation with mixed equipment (likely due to domain shift) suggests opportunities for domain adaptation in future work.

\section{Computational Efficiency}

\subsection{Inference Time}

Mean inference time per image on various hardware:

\begin{itemize}
    \item NVIDIA V100 GPU: 2.8 seconds (std 0.3s)
    \item NVIDIA T4 GPU: 3.5 seconds (std 0.4s)
    \item NVIDIA RTX A5000: 2.2 seconds (std 0.2s)
    \item CPU (Intel Xeon Gold 6248): 45 seconds (impractical for production)
\end{itemize}

GPU-accelerated inference enables clinical deployment with acceptable latency.

\subsection{Throughput and Scalability}

System throughput on single V100 GPU:

\begin{itemize}
    \item Single process: 300 images/hour
    \item Multiple processes (4): 950 images/hour
    \item Batch processing (batch size 32): 1,200 images/hour
\end{itemize}

Demonstrates ability to handle institutional volume requirements.

\section{Clinical Utility Analysis}

\subsection{Impact on Radiologist Decision-Making}

Testing with \underline{\hspace{2cm}} radiologists examining \underline{\hspace{2cm}} test images with and without LungInsight recommendations:

\begin{itemize}
    \item Diagnostic Accuracy Improvement: \underline{\hspace{0.5cm}}\% with system assistance
    \item Time to Diagnosis: \underline{\hspace{0.5cm}}\% reduction
    \item Radiologist Confidence: Increased by \underline{\hspace{0.5cm}}\% points (measured on 10-point scale)
    \item Agreement with System: \underline{\hspace{2cm}}\% on final diagnoses
\end{itemize}

Results suggest positive clinical impact though controlled trials needed for definitive evidence.

\subsection{False Positive and False Negative Analysis}

For clinical decision support, different error types have different implications:

\begin{itemize}
    \item \textbf{False Negatives (missed diseases):} More clinically concerning; mean system sensitivity 95.7\% is acceptable but not perfect
    \item \textbf{False Positives (over-diagnosis):} Can lead to unnecessary intervention; mean specificity 94.4\% is acceptable
\end{itemize}

Clinical protocols should account for model limitations, using system as adjunct rather than replacement for radiologist review.

\section{Limitations and Discussion}

\subsection{Dataset Limitations}

\begin{enumerate}
    \item Training data predominantly from developed countries; performance on global populations unknown
    \item Limited pediatric data; system primarily trained on adult chest X-rays
    \item Rare disease variants underrepresented in training data
    \item Dataset labels from automated systems and multiple annotators; potential label noise
\end{enumerate}

\subsection{Technical Limitations}

\begin{enumerate}
    \item Single imaging modality (2D chest X-rays); CT scans and other modalities not supported
    \item Image registration and temporal comparison not incorporated
    \item No integration of clinical context (patient history, age, comorbidities)
    \item Limited to binary classification per disease; severity assessment not provided
\end{enumerate}

\subsection{Clinical Limitations}

\begin{enumerate}
    \item System should assist, not replace, radiologist judgment
    \item Integration into clinical workflows requires institutional adaptation
    \item Regulatory approval and certification required for clinical deployment
    \item Potential medicolegal implications require careful institutional review
\end{enumerate}

\section{Discussion}

LungInsight demonstrates strong diagnostic performance across multiple lung diseases, achieving accuracy and sensitivity comparable to or exceeding published results for similar systems. The ensemble approach yielded consistent improvements over individual models, validating the multi-model strategy. Attention map visualizations enhanced interpretability and radiologist confidence in system recommendations.

The system shows promise for clinical deployment as a decision support tool, potentially reducing diagnostic errors and increasing radiologist efficiency. However, integration into actual clinical workflows requires addressing regulatory, operational, and ethical considerations beyond the technical scope of this project.

Performance variations across datasets and imaging equipment highlight the importance of ongoing validation and monitoring in deployed systems. Future work should focus on robustness across diverse populations and imaging protocols, integration of clinical context, and development of severity assessment capabilities.

}
